\section{Compact-sector Independence and Minimal Closure Requirements}
\label{sec:pr1-compact-countermodels}

The seven PR-I postulates establish the common geometric foundation.  The compact construction is completed by the PR-I boundary-orientation and spectral-minimality principles.  This section formalizes those sector-specific conditions and proves that they select the released compact equations without a diagnostic input.  We hold fixed the master field, internal cylinder, metric, Hilbert space, separable spectral operator, projection-trace radial potential, trace alphabet $T_A=1/4$, $T_X=3/4$, and the primitive-return denominator
\begin{equation}
 D_A=1-T_A^3-T_A^4=\frac{251}{256}.
\end{equation}
No diagnostic input is admitted.

For $\sigma\in\{-1,0,+1\}$ define the finite, integer-coefficient terminal
cofactor
\begin{equation}
 N_\sigma=1+T_A^3+T_A^4+T_X\left(T_A^5+\sigma T_A^8\right).
 \label{eq:pr35-terminal-countermodels}
\end{equation}
Each member contains no continuous free parameter.  Nevertheless the three members yield distinct rational values of $N_\sigma$, $R_\sigma=N_\sigma/D_A$,
and
\begin{equation}
 c_\sigma=T_X\left(1+T_A^2-T_A^6R_\sigma\right).
\end{equation}
The canonical PR coefficient is $\sigma=-1$.  The pre-projection exclusion equation has exact defect $|1+\sigma|$, so $\sigma=-1$ is its sole zero; the other displayed coefficients are comparison values outside the completed PR-I boundary class. For exhaustive comparison in the nondegenerate radial eigenbasis, consider
\begin{equation}
 \Pi_{1j}=|\phi_1\rangle\langle\phi_1|+
 |\phi_j\rangle\langle\phi_j|,\qquad j\in\{3,5,7\},
\end{equation}
Each is a rank-two, parity-preserving orthogonal spectral projector with no continuous coefficient.  The projector comparison is resolved by the required-channel minimality rule.  Among coordinate projectors containing the fundamental winding-linked channel $\phi_1$ and primitive spatial-return
channel $\phi_3$, minimal rank uniquely gives $\Pi_{13}$.  Projectors containing different odd channels do not satisfy that declared PR-I closure condition.

Thus the signed-incidence and required-channel minimality certificates each have one admissible PR orbit.  Their Cartesian certificate classifies all nine displayed coefficient/projector combinations and retains only $(-1,\Pi_{13})$.  These target-independent closure principles are logical
extensions of the PR foundation and need not be literal restatements of Postulates 1 throught 7.
